% Sections 13.1 to 13.6, translated from sv/chapters/13/theory.tex.

\section{Derivatives of trigonometric functions}
\label{c13:sec:trig}

\begin{narrowtable}{@{}ll@{}}
  \thead{Function} & \thead{Derivative} \\
  \midrule
  $\sin(x)$ & $\cos(x)$ \\
  $\cos(x)$ & $-\sin(x)$ \\
  $\sin(kx)$ & $k\cos(kx)$ \\
  $\cos(kx)$ & $-k\sin(kx)$ \\
\end{narrowtable}

\begin{exempel}
\[
  f(x) = 3\sin(2x) \quad \Rightarrow \quad f'(x) = 3 \cdot 2\cos(2x) = 6\cos(2x)
\]
\end{exempel}

\section{Derivatives of exponential functions}
\label{c13:sec:exp}

\begin{narrowtable}{@{}ll@{}}
  \thead{Function} & \thead{Derivative} \\
  \midrule
  $e^x$ & $e^x$ \\
  $e^{kx}$ & $ke^{kx}$ \\
  $a^x$ & $a^x \ln a$ \\
\end{narrowtable}

\textbf{Derivation of $\frac{d}{dx}e^x = e^x$.} With the definition of the derivative from
\secref{c12:sec:definition},
\[
  \lim_{h \to 0} \frac{e^{x+h} - e^x}{h} = e^x \lim_{h \to 0} \frac{e^h - 1}{h} = e^x \cdot 1
  = e^x.
\]

\section{Derivatives of logarithmic functions}
\label{c13:sec:log}

{\renewcommand{\arraystretch}{2}%
\begin{narrowtable}{@{}ll@{}}
  \thead{Function} & \thead{Derivative} \\
  \midrule
  $\ln(x)$ & $\dfrac{1}{x}$ \\
  $\ln(kx)$ & $\dfrac{1}{x}$ \\
  $\log_{10}(x)$ & $\dfrac{1}{x \ln 10}$ \\
\end{narrowtable}}

\begin{aside}
\note[Note] $\ln(kx) = \ln k + \ln x$, whose derivative is $1/x$: the constant $\ln k$ disappears.
\end{aside}

\section{The product rule and the chain rule}
\label{c13:sec:regler}
The tables above give the derivative of a single function. When two functions are multiplied
together, or when one function is substituted into another, two more rules are needed.

\subsection{The product rule}
\label{c13:sec:produktregeln}
The derivative of a product $f(x) = g(x) \cdot h(x)$ is
\[
  f'(x) = g'(x) \cdot h(x) + g(x) \cdot h'(x).
\]
So differentiate one factor at a time while the other stays as it is, and add the two terms.

\begin{exempel}
$f(x) = x^2 \sin(x)$ with $g(x) = x^2$ and $h(x) = \sin(x)$ gives
\[
  f'(x) = 2x \cdot \sin(x) + x^2 \cdot \cos(x).
\]
\end{exempel}

\subsection{The chain rule}
\label{c13:sec:kedjeregeln}
A composite function $f(x) = g(u(x))$ consists of an \textbf{outer function} $g$ and an
\textbf{inner function} $u(x)$. Its derivative is the derivative of the outer function, evaluated
at the inner function, times the derivative of the inner function:
\[
  f'(x) = g'(u(x)) \cdot u'(x).
\]

\begin{exempel}
$f(x) = \ln(2x + 1)$ has the outer function $\ln u$ and the inner function $u = 2x + 1$, with
$u' = 2$:
\[
  f'(x) = \frac{1}{2x + 1} \cdot 2 = \frac{2}{2x + 1}.
\]
\end{exempel}

The table rows for $\sin(kx)$, $\cos(kx)$, $e^{kx}$ and $\ln(kx)$ above are special cases of the
chain rule, with the inner function $u = kx$ and $u' = k$. For example,
\[
  \frac{d}{dx}\sin(kx) = \cos(kx) \cdot k = k\cos(kx), \qquad
  \frac{d}{dx}\ln(kx) = \frac{1}{kx} \cdot k = \frac{1}{x}.
\]

\section{Table of derivatives}
\label{c13:sec:tabell}

\begin{narrowtable}{@{}ll@{}}
  \thead{Function $f(x)$} & \thead{Derivative $f'(x)$} \\
  \midrule
  $x^n$ & $nx^{n-1}$ \\
  $e^x$ & $e^x$ \\
  $e^{kx}$ & $ke^{kx}$ \\
  $\ln x$ & $1/x$ \\
  $\sin x$ & $\cos x$ \\
  $\cos x$ & $-\sin x$ \\
  $\sin(kx)$ & $k\cos(kx)$ \\
  $\cos(kx)$ & $-k\sin(kx)$ \\
\end{narrowtable}

\section{Worked examples}
\label{c13:sec:typexempel}

\begin{exempel}[Derivatives of common functions]
Differentiate \textbf{a)} $f(x) = 5e^{3x}$, \textbf{b)} $f(x) = \ln(4x)$ and
\textbf{c)} $f(x) = 2\cos(x) - 3\sin(x)$.
\begin{align*}
  \textbf{a)}\quad f'(x) &= 5 \cdot 3e^{3x} = 15e^{3x} \\
  \textbf{b)}\quad f'(x) &= 1/x \\
  \textbf{c)}\quad f'(x) &= -2\sin(x) - 3\cos(x)
\end{align*}
\end{exempel}

\begin{exempel}[The derivative at a point]
The current $i(t) = I_0 e^{-t/\tau}$ amperes describes the discharge of a capacitor. Calculate
$i'(0)$.
\[
  i'(t) = -\frac{I_0}{\tau} e^{-t/\tau} \quad \Rightarrow \quad i'(0) = -\frac{I_0}{\tau}
\]
$i'(0)$ is the instantaneous rate of change of the current at $t = 0$.
\end{exempel}

\begin{exempel}[Stationary point of an exponential function]
Find the turning point of $f(x) = xe^{-x}$. The product rule (\secref{c13:sec:produktregeln})
with $g(x) = x$ and $h(x) = e^{-x}$ gives
\[
  f'(x) = e^{-x} + x \cdot (-e^{-x}) = e^{-x}(1 - x)
\]
$f'(x) = 0$ gives $x = 1$, since $e^{-x} \neq 0$. The second derivative is
\[
  f''(x) = -e^{-x}(1 - x) + e^{-x}(-1) = e^{-x}(x - 2),
\]
and
\[
  f''(1) = e^{-1}(1 - 2) = -e^{-1} < 0 \quad \Rightarrow \quad
  \text{\textbf{maximum} at } (1, e^{-1}) \approx (1, 0.37).
\]
\end{exempel}
